% !TeX program = xelatex
\documentclass[11pt,a4paper,UTF8]{ctexrep}
\usepackage{geometry}
\geometry{left=2.5cm,right=2.5cm,top=2.5cm,bottom=2.5cm}
\usepackage{amsmath,amssymb,amsfonts,amsthm}
\usepackage{booktabs}
\usepackage{graphicx}
\usepackage{hyperref}
\usepackage{tcolorbox}
\usepackage{titlesec}
\usepackage{fancyhdr}
\usepackage{listings}
\usepackage{xcolor}

\pagestyle{fancy}
\fancyhf{}
\fancyhead[L]{\small\kaishu 数学建模期末冲刺复习讲义}
\fancyhead[R]{\small\thepage}
\fancyfoot[C]{\thepage}

\tcbset{
    sharp corners,
    colback=white,
    boxrule=0.8pt,
    leftrule=3pt
}

\newtcolorbox{coretheorem}[1]{
    colframe=blue!70!black,
    title=\textbf{#1},
    coltitle=white,
    colbacktitle=blue!70!black
}

\newtcolorbox{modelbox}[1]{
    colframe=teal!70!black,
    title=\textbf{#1},
    coltitle=white,
    colbacktitle=teal!70!black
}

\newtcolorbox{pitfallbox}[1]{
    colframe=red!70!black,
    title=\textbf{#1},
    coltitle=white,
    colbacktitle=red!70!black
}

\title{\Huge\textbf{数学建模期末复习与考点精讲手册}\\[0.5cm]
\Large 基于祁永强《数学建模》与课件深度提炼}
\author{\textbf{授课助教 / 讲义编写组}}
\date{\today}

\begin{document}
\maketitle
\tableofcontents
\newpage

\chapter{经典初等与微分动力学模型}

\section{机理分析与连续微积分模型}

\begin{modelbox}{核心模型：双层玻璃窗保温机理}
设两层玻璃厚度均为 $d$，导热系数为 $k_1$；中间静止空气层厚度为 $l$，导热系数为 $k_2$。
由一维稳态热传导傅里叶定律：
\begin{equation}
    Q = \frac{\Delta T}{R_{\text{total}}}, \quad R_{\text{total}} = \frac{2d}{k_1} + \frac{l}{k_2}
\end{equation}
因空气导热系数极小（$k_1/k_2 \approx 16 \sim 32$），中间空气层贡献了绝大部分热阻，流失热量较单层玻璃大幅减少。
\end{modelbox}

\begin{pitfallbox}{易错警示：微分方程相平面分析陷阱}
在分析传染病 SIR 模型或兰彻斯特平方律时，切忌直接对时间变量 $t$ 积分求显式解（通常是超越方程或不可积）。应当通过消去时间参量微元：
\begin{equation}
    \frac{\mathrm{d}i}{\mathrm{d}s} = \frac{\beta s i - \gamma i}{-\beta s i} = -1 + \frac{\gamma}{\beta s}
\end{equation}
求出第一积分（相轨线解析关系式），再结合初始边界值确定临界阈值与渐近行为。
\end{pitfallbox}

\chapter{运筹规划与多目标决策}

\section{层次分析法（AHP）答题标准范式}

\begin{coretheorem}{AHP 考点清单与一致性检验判据}
\begin{enumerate}
    \item \textbf{判断矩阵构造}：$A = (a_{ij})_{n \times n}$ 满足 $a_{ji} = 1/a_{ij}, a_{ii} = 1$。
    \item \textbf{最大特征根与特征向量近似计算}：和法或根法归一化。
    \item \textbf{一致性指标}：$CI = \frac{\lambda_{\max} - n}{n - 1}$。
    \item \textbf{一致性比率}：$CR = \frac{CI}{RI}$。当 $CR < 0.10$ 时接受一致性，否则必须调整矩阵。
\end{enumerate}
\end{coretheorem}

\end{document}
